% Avsnitt 18.1-18.8, ur lectures/L18/appendix/a_register_bank.md.

\section{Varför det här lagret finns}
\label{c18:sec:why}

\code{can_controller}s statusutgångar är gjorda för en anropare i samma klockdomän som ser varje
cykel: \code{tx_done} och \code{rx_valid} är \textbf{encykelspulser}, 20~ns vid 50~MHz
(\chapref{c11:ch}). En drivrutin som pollar över SPI samplar systemet några \emph{mikrosekunder} i
taget i bästa fall \textemdash{} en puls är borta tusentals cykler innan nästa pollning kommer.
Registerkartan (bilaga~\ref{app:protocol}) lovar därför något annat: \textbf{klibbiga, pollbara
nivåer} med uttryckliga nollställningar (\code{RX_ACK}, \code{ERROR_FLAGS}), och skrivutlösta
händelser (\code{TX_SEND}).

Att konvertera mellan de två världarna är hela den här modulens uppgift. Den innehåller ingen
protokollkunskap i någon riktning: ingenting om CAN-ramar (det stannar i \code{can_controller}) och
ingenting om SPI (det kommer senare i det här kapitlet, och i \chapref{c19:ch}). Det är register,
låsningar och en pulsgenerator \textemdash{} vilket är precis därför den kan verifieras fullständigt
under ett enda pass, innan det finns någon SPI alls.

Du har sett båda dess kärntrick förut, i \chapref{c3:ch} och \ref{c4:ch}. En klibbig statusbit är
\chapref{c3:ch}:s vippa med de två frågorna uttryckligen besvarade: vad som sätter den, och vad som
nollställer den. Och att göra om en hållen nivå till en encykelspuls är precis vad
\code{button_sync}s flankdetektering gjorde i \chapref{c4:ch} \textemdash{} här går den åt samma håll
för \code{TX_SEND} (en skrivning blir en puls) och åt motsatt håll för \code{tx_done} (en puls blir
en nivå som ligger kvar). Det här kapitlet befordrar båda från en övning till ett riktigt
arkitekturlager med ett dokumenterat kontrakt.

% ------------------------------------------------------------------------------------------
\section{Gränssnittet}
\label{c18:sec:iface}

Sexton portar, \textbf{alla ingångar först, sedan alla utgångar}, samma konvention som varje modul i
projektet \textemdash{} och, som överallt annars i det, är bindningen \emph{positionell}, så
ordningen nedan är det kontrakt \code{register_bank_tb} håller dig till:

\begin{textdrawing}
 1 clock        5 reg_wdata    9 rx_data     13 tx_id
 2 reset_s2_n   6 tx_done     10 rx_valid    14 tx_dlc
 3 reg_index    7 rx_id       11 error       15 tx_data
 4 reg_write    8 rx_dlc      12 reg_rdata   16 tx_req
\end{textdrawing}

Typerna kommer från \code{can_def} (\code{use work.can_def.all;}), så portarna mot kontrollern har
exakt samma typer som \code{can_controller}s egna. Vid sidan av \code{clock} och \code{reset_s2_n}
(portarna 1 och 2) faller de fjorton övriga i två grupper.

\paragraph{Sidan mot bryggan (portarna 3--5 och 12)}
En medvetet minimal registeråtkomstport \textemdash{} det som \chapref{c19:ch}:s brygga kommer att
driva, och det som testbänken driver direkt i dag:
\begin{itemize}
  \item \code{reg_index : std_logic_vector(3 downto 0)} \textemdash{} vilket register, protokollets
    index (offset delat med 4).
  \item \code{reg_write : std_logic} \textemdash{} en \textbf{encykelsstrob}; på dess stigande
    klockflank verkställs \code{reg_wdata} i det adresserade registret.
  \item \code{reg_wdata : std_logic_vector(31 downto 0)} \textemdash{} värdet som ska skrivas.
  \item \code{reg_rdata : std_logic_vector(31 downto 0)} \textemdash{} det adresserade registrets
    aktuella värde, \textbf{kombinatoriskt}: det följer \code{reg_index} utan att någon klockflank är
    inblandad. Regeln om att låsa en gång vid läsning ur protokollspecifikationen är \emph{bryggans}
    uppgift (\chapref{c19:ch}), inte den här modulens \textemdash{} banken berättar alltid den
    nuvarande sanningen, och bryggan avgör när den tar sitt enda atomära sampel.
\end{itemize}

\paragraph{Sidan mot kontrollern (portarna 6--11 och 13--16)}
Speglar \code{can_controller}s portar på registersidan en mot en: \code{tx_id}, \code{tx_dlc},
\code{tx_data} och \code{tx_req} ut till kontrollern, och \code{tx_done}, \code{rx_id},
\code{rx_dlc}, \code{rx_data}, \code{rx_valid} och \code{error} tillbaka från den. I
\chapref{c19:ch}:s toppnivå kopplas de rakt över, port till port, ingenting emellan.

\code{reset_s2_n} är den \emph{synkroniserade} reseten, som för varje delblock i projektet.
\Chapref{c19:ch}:s toppnivå äger en \code{meta_prev} för den, precis som \code{can_controller} äger
sin egen internt.

% ------------------------------------------------------------------------------------------
\section{De tre STATUS-bitarna: en låsning var}
\label{c18:sec:status}

Varje \code{STATUS}-bit är en vippa med ett namngivet sättvillkor och ett namngivet
nollställningsvillkor, och ingenting mer:

\begin{booktable}{@{}p{2.2em}p{5em}L>{\raggedright\arraybackslash}p{9em}@{}}
  \thead{Bit} & \thead{Betydelse} & \thead{Sätts av} & \thead{Nollställs av} \\
  \midrule
  0 & TX klar & reset; kontrollerns puls \code{tx_done}; \textbf{en stigande flank på \code{error}
    medan den här biten är låg}; \textbf{en accepterad skrivning till \code{TX_ABORT}} & en
    accepterad skrivning till \code{TX_SEND} \\
  1 & RX giltig & kontrollerns puls \code{rx_valid} & en skrivning av \code{0x1} till
    \code{RX_ACK} \\
  2 & Fel & en \textbf{stigande flank} på kontrollerns nivå \code{error} & en skrivning av
    \code{0x0} till \code{ERROR_FLAGS} \\
\end{booktable}

Tre av dem förtjänar en andra blick:
\begin{itemize}
  \item \textbf{Bit 0 börjar på \code{'1'}.} Ur reset är kontrollern i vila, så banken är redo att ta
    emot en ram \textemdash{} en drivrutin som pollar innan den någonsin sänt måste se ”klar”, annars
    skulle \code{send()} misslyckas för alltid på ett nyss spänningssatt system.
  \item \textbf{Bit 2 låser en flank, inte en nivå.} Kontrollern håller \code{error} hög tills nästa
    ram börjar (\chapref{c11:ch}), på sitt eget schema. Att låsa den stigande flanken ger
    \emph{drivrutinen} en egen livscykel för flaggan: en nollställning tar bort den även medan
    kontrollerns nivå fortfarande är hög, och låsningen sätts inte igen förrän ett genuint \emph{nytt}
    fel producerar en ny stigande flank. Att i stället låsa på nivån skulle få nollställningen att
    verka trasig \textemdash{} nollställd och omedelbart satt igen \textemdash{} så länge kontrollern
    håller ledningen.
  \item \textbf{Bit 0 har tre vägar tillbaka upp, och det är avsiktligt.} Läs de två extra som
    livrem och hängslen kring det enda fel som kan slå ut hela noden.
\end{itemize}

\begin{aside}
En sändning som förlorar arbitreringen slutar utan att någonsin nå EOF. \Chapref{c17:ch} låter
kontrollern pulsa \code{tx_done} också på den vägen, just för att den här biten ska komma tillbaka
\textemdash{} men kontrollern är \emph{din}, och en version som glömmer det lämnar biten låg för
alltid, så varje senare \code{TX_SEND} kastas och noden kan inte sända igen förrän den resettas.
Banken förlitar sig därför inte på att kontrollern gör rätt. \textbf{En stigande flank på
\code{error} medan bit 0 redan är låg kan bara betyda att en sändning slutade illa}, eftersom en nod
antingen sänder eller tar emot, och ett mottagningsfel inte kan inträffa medan en sändning pågår. Så
den flanken sätter biten också, och det behövs ingen \code{role}-port för att veta det.

\code{TX_ABORT} är den tredje vägen: en uttrycklig nödutgång som drivrutinen kan dra i om den bedömer
att den väntat länge nog. Ingenting i konstruktionen ska någonsin behöva den. Bygg in den ändå
\textemdash{} kostnaden är ett enda register som bara går att skriva till, och alternativet är en nod
som en drivrutin inte har någon möjlighet att få loss.
\end{aside}

En enda ordningsregel binder ihop låsningarna: i den klockade processen \textbf{tillämpas
kontrollerns händelser efter bryggans skrivningar}, så om en \code{TX_SEND}-skrivning och en
\code{tx_done}-puls någonsin landade på samma flank vinner hårdvaran och banken förblir klar. En så
sällsynt konflikt bör lösas åt det håll som inte kan köra fast systemet, och det är samma princip som
stycket ovan tillämpar tre gånger om.

% ------------------------------------------------------------------------------------------
\section{\code{TX_SEND}: en skrivning blir en puls}
\label{c18:sec:txsend}

En \code{TX_SEND}-skrivning med \code{reg_wdata(0) = '1'}, som anländer medan TX klar är satt, gör
exakt två saker på den klockflanken: höjer \code{tx_req} \textbf{under en cykel}, och nollställer TX
klar. Mer precist: \code{tx_req} är hög under den enda cykeln \emph{efter} skrivningens klockflank,
och låg igen på nästa \textemdash{} testbänken kontrollerar båda ändarna av det påståendet.

Allt annat med det här registret är vad det \emph{ignorerar}:
\begin{itemize}
  \item En skrivning medan TX klar är låg (en sändning pågår) kastas helt \textemdash{} ingen puls,
    ingen tillståndsändring. Att skriva över \code{tx_id}, \code{tx_dlc} eller \code{tx_data} mitt i
    en sändning är den korruption som drivrutinens egen klar-kontroll skyddar mot från
    mjukvarusidan; det här är samma skydd på hårdvarusidan, och ingendera sidan får anta att den
    andra finns där.
  \item En skrivning med \code{reg_wdata(0) = '0'} gör ingenting: registerkartan säger ”skriv
    \code{0x1} för att utlösa”, och banken håller den till det.
  \item Läsningar ger noll \textemdash{} \code{TX_SEND} lagrar ingenting, det \emph{gör} något.
\end{itemize}

Varför pulsdisciplinen spelar roll syns tydligast i felet den förhindrar. \code{can_controller} agerar
på \code{tx_req} så fort den är hög och bussen är ledig (\chapref{c16:ch}), så en \code{tx_req} som
lämnas \emph{hållen} begär på nytt i samma ögonblick som en ram tar slut, för alltid, och suddar
\code{error} vid varje passage tillbaka genom starttillståndet. Begäran måste vara en händelse, och
något måste göra den till en. På det här lagret är det enklare \textemdash{} en skrivstrob är redan
en händelse, inte en nivå \textemdash{} men det måste ändå \emph{upprätthållas}: en skrivning, en
puls, oavsett vad mastern gör härnäst.

% ------------------------------------------------------------------------------------------
\section{\code{TX_ABORT}: nödutgången}
\label{c18:sec:txabort}

Index 12, bara skrivbart, och spegelbilden av \code{TX_SEND}: en skrivning av \code{0x1} sätter TX
klar tillbaka till \code{'1'}. En skrivning av något annat gör ingenting, och en läsning ger
\code{0x00000000}.

Det rör ingenting annat. Det når inte \code{can_controller} \textemdash{} kontrollern har ingen
ingång för att avbryta, dess femton portar spikades i \chapref{c11:ch} och \code{can_controller_tb}
binder dem positionellt \textemdash{} och den behöver ingen. En kontroller som avbrutit en sändning
är redan tillbaka i \code{STATE_IDLE} och fullt kapabel att sända igen; det enda som sitter fast är
den här bankens klibbiga bit. \code{TX_ABORT} lossar den, och det är hela modulen.

Det är värt att stanna vid ett ögonblick, eftersom det är buggens allmänna form. Ingenting i
\emph{hårdvaran} hade kört fast. Det som kört fast var en \textbf{statusbit vars enda sättvillkor
hängde på en händelse som inte längre inträffade}. Varje klibbig flagga med en väg upp och en väg ner
är en missad händelse från samma fel, och lösningen är alltid densamma: se till att varje väg ut ur
upptagettillståndet också höjer flaggan, och ge anroparen ett sätt att tvinga fram den ifall du har
fel om ”varje”.

% ------------------------------------------------------------------------------------------
\section{RX-sidan: infångning vid pulsen}
\label{c18:sec:rx}

När \code{rx_valid} pulsar gör banken två saker på samma flank: sätter STATUS bit 1, och
\textbf{fångar in} \code{rx_id}, \code{rx_dlc} och \code{rx_data} i de fyra \code{RX_*}-registren
\textemdash{} \code{rx_data}s övre ord (bitarna 63--32, databyte 0--3, byte 0 i den mest signifikanta
positionen) i \code{RX_DATA_LO}, dess undre ord i \code{RX_DATA_HI}.

Att fånga in, i stället för att koppla kontrollerns portar rakt till läsmuxen, är det som ger
drivrutinen en stabil ram att läsa: \code{RX_*}-registren håller sina värden medan drivrutinen
arbetar sig igenom fyra separata SPI-lästransaktioner, oavsett vad kontrollerns portar gör under
tiden. En skrivning till \code{RX_ACK} nollställer \emph{bara flaggan} \textemdash{} dataregistren
behåller sitt innehåll till nästa infångning, vilket är ofarligt: registerkartans kontrakt är att
\code{RX_*} är meningsfullt \emph{medan RX giltig är satt}.

Vad infångningen \textbf{inte} ger drivrutinen är buffring. En ny ram som fullbordas innan den förra
kvitterats fångas helt enkelt in ovanpå den \textemdash{} en ram hålls, den nyaste vinner. Det är
inte den här modulen som är lat; det är en dokumenterad lucka som ärvts från kontrollern själv, som
håller exakt en mottagen ram och skriver över den vid nästa (\chapref{c11:ch}). Banken gör den luckan
synlig på registernivå i stället för att dölja den, och \chapref{c19:ch} återkommer till vad den
kostar en drivrutin.

% ------------------------------------------------------------------------------------------
\section{Maskning vid skrivning}
\label{c18:sec:mask}

Skrivbara register lagrar bara de bitar som finns i hårdvaran, och läses tillbaka som nollor ovanför
dem:

\begin{booktable}{@{}p{9em}p{9em}L@{}}
  \thead{Register} & \thead{Lagrade bitar} & \thead{En skrivning av \code{0xFFFFFFFF} läses tillbaka
    som} \\
  \midrule
  \code{TX_ID} & 10--0 (\code{id_t}) & \code{0x000007FF} \\
  \code{TX_DLC} & 3--0 (\code{dlc_t}) & \code{0x0000000F} \\
  \code{TX_DATA_LO}/\code{HI} & alla 32 & \code{0xFFFFFFFF} \\
\end{booktable}

Lägg märke till vad banken \emph{inte} gör: validerar. En \code{TX_DLC} på \code{0xF} lagras och
lämnas vidare till kontrollern som den är \textemdash{} att avvisa \code{dlc > 8} är drivrutinens
ansvar, i den parallella mjukvaruklassen, och att göra det två gånger på två språk är precis så de
två kopiorna glider isär. Banken maskar till de bitar som fysiskt finns; drivrutinen upprätthåller
betydelsen. \code{STATUS} och \code{RX_*}-registren ignorerar skrivningar helt, och de reserverade
indexen 13--15 läses som noll och ignorerar skrivningar, precis som protokollspecifikationen kräver.

Utgångarna mot kontrollern kommer direkt från TX-registren:
\code{tx_data <= TX_DATA_LO & TX_DATA_HI} \textemdash{} 64 bitar med databyte 0 i den mest
signifikanta positionen, den layout \code{can_controller} väntar sig (\chapref{c16:ch}) och samma
packning som drivrutinen utför på andra sidan kartan (se diagrammet över databyteordningen i
bilaga~\ref{app:protocol}).

% ------------------------------------------------------------------------------------------
\section{Att skriva \code{register_bank.vhd}}
\label{c18:sec:write}

Filen hamnar i \file{bridge/}, bredvid dina kopierade kontrollermoduler. Formen är två processer och
en handfull konkurrenta tilldelningar:
\begin{itemize}
  \item \textbf{En klockad process} som äger varje register och låsning: resetvärdena (TX klar
    \code{'1'}, allt annat noll), skriv-\code{case}:et över \code{reg_index}, och sedan de tre
    uppdateringarna från kontrollerns händelser, i den ordningen.
  \item \textbf{En kombinatorisk process} för läsmuxen: ge \code{reg_rdata} förvalet idel nollor, och
    skriv sedan över de bitar varje register faktiskt har \textemdash{} samma stil med ”tilldela ett
    förval, skriv över i specifika grenar” som varje räknare och tillståndsmaskin från
    \chapref{c6:ch} och framåt.
  \item Konkurrenta tilldelningar som kopplar ut \code{tx_id}, \code{tx_dlc}, \code{tx_data} och
    \code{tx_req} från TX-registren och pulsvippan.
\end{itemize}

Analysera den och kör sedan den utdelade testbänken. \file{can_def.vhd} ligger en katalog bort, i
\file{controller/}: banken läser \code{can_def} för sina typer mot kontrollern men hör hemma på
SPI-sidan av protokollgränsen.

\begin{shell}
cd bridge
ghdl -a --std=93 ../controller/can_def.vhd register_bank.vhd register_bank_tb.vhd
ghdl -e --std=93 register_bank_tb
ghdl -r --std=93 register_bank_tb --assert-level=error
\end{shell}

\paragraph{Vad testbänken kontrollerar}
\file{register_bank_tb.vhd} går igenom kontraktet i elva fall, i tur och ordning: reset-tillståndet
(bit 0 satt, allt annat noll); skrivning och återläsning med maskning på varje skrivbart register,
och utgångarna mot kontrollern som följer dem; \code{TX_SEND}s encykelspuls och nollställningen av
klar; skrivningen som kastas medan sändning pågår; \code{tx_done} som återställer klar; en
\code{TX_SEND}-skrivning av \code{0x0} som inte gör något; infångning av mottagen ram, att den håller
trots att ingångarna ändras, att den är skrivskyddad, och \code{RX_ACK}; fellåsningens flanksemantik,
inklusive regeln om skrivning av noll och en nollställning medan nivån fortfarande är hög; de
reserverade indexen; och sist de två återhämtningsvägar som ger tillbaka TX klar efter en förlorad
arbitrering, en automatisk och en driven av \code{TX_ABORT}.

Varje fallkommentar i filen säger vad ett underkänt fall betyder, så läs testbänken innan du skriver
modulen \textemdash{} den är den körbara formen av det här kapitlet. De två sista fallen är särskilt
värda att läsa i förväg: utan dem fastnar noden efter sin första förlorade arbitrering, och eftersom
en tvånodsbuss provocerar det i normal drift är det inte ett hörnfall.
